\begin{dang*}{Thay đổi số vòng dây máy biến áp}
%\paragraph{Bài toán liên quan đến số vòng dây của máy biến áp thay đổi}
%\Binhluan{
%	\begin{itemize}
%		\item Khi máy biến áp có số vòng dây ở cuộn sơ cấp thay đổi: 
%		\begin{align*}
%		\boder{\dfrac{U_2}{U_1} = \dfrac{N_2}{N_1};\quad 
%		\dfrac{U'_2}{U_1} = \dfrac{N_2}{N_1\pm n}}
%		\end{align*}
%		\item Khi máy biến áp có số vòng dây ở cuộn thứ cấp thay đổi: 
%		\begin{align*}
%		\boder{\dfrac{U_2}{U_1} = \dfrac{N_2}{N_1};\quad
%		\dfrac{U'_2}{U_1} = \dfrac{N_2\pm n}{N_1}}
%		\end{align*}
%	\end{itemize} 
% }
%{
%Mấu chốt nằm ở bài toán cho điện áp đặt vào hai đầu cuộn sơ cấp không đổi
%\begin{align*}
%	U_1=const
%\end{align*}
%}
%\Binhluan{
%\paragraph{Bài toán máy biến áp có một số vòng dây quấn ngược}
%Nếu một cuộn dây nào đó (ví dụ cuộn sơ cấp) có n vòng dây quấn ngược thì từ trường của n vòng này ngược với từ trường của phần còn lại nên nó có tác dụng khử bớt từ trường của n vòng dây còn lại.
%\begin{align*}
%\boder{\dfrac{U_2}{U_1} = \dfrac{N_2}{N_1-2n}}
%\end{align*}
%}
%{
%Ta coi như cuộn dây này bị mất đi 2n vòng.
%}
\end{dang*}

